\par\Needspace{6\baselineskip}\medskip
\hypertarget{native-142-article103}{}\pdfbookmark[2]{103.}{article-103}
\begin{center}103.\end{center}
Whenever $z$ takes a larger value beyond the limits of Table III, equations 13 and $13^*$ can always be solved safely and conveniently by trial in their original form, for reasons similar to those we explained for the ellipse in Art. 94. In such a case one may \SourceLeaf{143}assume that the orbital elements are known at least roughly. An approximate value of $n$ is then obtained at once from the formula $\operatorname{tang}2n=\dfrac{\sin f\mathpunct{.}\BareRoot{rr'}}{\alpha\BareRoot{(ee-1)}}$, which follows directly from equation 6 of Art. 99. From $n$ we obtain $z$ by the formula $z=\dfrac{1-\cos2n}{2\cos2n}=\dfrac{\sin n^2}{\cos2n}$. Starting with the approximate value of $z$, a few trials give the value that satisfies equation 13 or $13^*$ exactly. Those equations can also be written in the form
\[
\begin{aligned}
m={}&\left(l-\frac{\sin n^2}{\cos2n}\right)^{\frac12}\\[.6em]
&+2\left(l-\frac{\sin n^2}{\cos2n}\right)^{\frac32}\frac{\dfrac{\operatorname{tang}2n}{\cos2n}-\log\operatorname{hyp}\operatorname{tang}(45^0+n)}{\operatorname{tang}2n^3}
\end{aligned}
\]
\[
\begin{aligned}
M={}&-\left(L+\frac{\sin n^2}{\cos2n}\right)^{\frac12}\\[.6em]
&+2\left(L+\frac{\sin n^2}{\cos2n}\right)^{\frac32}\frac{\dfrac{\operatorname{tang}2n}{\cos2n}-\log\operatorname{hyp}\operatorname{tang}(45^0+n)}{\operatorname{tang}2n^3}
\end{aligned}
\]
\noindent and in this way, by dispensing with $z$, we can find the true value of $n$ directly.
